\vspace{-0.25cm}
\section{Related Work}
\vspace{-0.25cm}
\label{sec:related}


\mysection{Communicative Agents} 
Communication between agents has been studied for a long time \cite{minsky1988society,minsky2007emotion}. There are many ways to facilitate communication between agents, and with agents \cite{finin1994kqml,poslad2007specifying,russell2010artificial}. Among these, natural language is considered the most natural form of communication \cite{russell2010artificial}.
By enabling agents to function as communicators themselves, they become capable of solving complex tasks \cite{Tan1997MultiAgentRL, Panait2005CooperativeML, lowe2017multi,andreas2022language,foerster2016learning,sukhbaatar2016learning,mordatch2018emergence,havrylov2017emergence,du2021learning,sheng2022learning,nakano2021webgpt,shuster2022blenderbot,amelia2022improving,hosseini2020simple,ahn2022do,huang2022inner,abramson2020imitating,siddaharth2021lila,huang2022language,li2022pretrained,reed2022gato}. Communication between AI agents can occur in a competitive setting \cite{tesauro1995temporal, silver2017mastering} or a cooperative setting \cite{hadfield2016cooperative,dafoe2020open,bard2020hanabi,zhu2023chatgpt,lo2023cheap}. Cooperative AI refers to artificial intelligence systems that are designed to work together with humans and other AI systems to achieve common goals \cite{Claus1998TheDO, wooldridge2009introduction}. Cooperative AI systems take into account the needs and capabilities of other agents in the system and actively seek to collaborate and coordinate their actions with them, which has many potential benefits, including increased efficiency, improved decision-making, and the ability to tackle complex problems that are beyond the reach of any single agent. However, designing effective cooperative AI systems is still an active area of research, as it requires addressing a range of technical, ethical, and social challenges \cite{dafoe2020open}. Our work enables communicative agents to engage in a conversation and cooperate with each other to solve assigned tasks. The agents, each assigned a distinct role, are expected to apply their expertise and knowledge to solve their common task.



\mysection{Instructional LLMs and Prompt Engineering} LLMs are trained on diverse text data and excel in text completion, with various downstream NLP applications \cite{brown2020language,chowdhery2022palm,hoffmann2022chinchilla,zhang2022opt,touvron2023llama}. However, InstructGPT suggests that LLMs may not align with user intent, proposing reinforcement learning from human feedback (RLHF) \cite{christiano2017deep} and Instruction Fine-Tuning (IFT) \cite{wei2021flan} to improve LLMs' relevance and appropriateness to user instructions.
Special types of instruction or prompting methods , such as Chain-of-Thought (CoT) \cite{wei2022chain}, zero-shot-CoT \cite{kojima2022large}, and ReAct \cite{yao2023react}, have recently been developed to enhance the performance of LLMs on reasoning, arithmetic and decision making tasks \cite{zhou2022least,wang2023self,imani2023mathprompter,lu2023dynamic,fu2022complexity,shi2022language,hendrycks2021measuring,lewkowycz2022solving,zhang2023automatic,ho2022large,zhang2023multimodal,shinn2023reflexion,zelikman2022star,creswell2022selection,nye2021show,sridhar2023hierarchical}.
These techniques underpin the impressive capabilities of recent dialogue LLMs \cite{shuster2022blenderbot,thoppilan2022lamda,glaese2022improving,bai2022training,openai_chatgpt,bubeck2023sparks}, which aim to simulate human-like conversations and provide personalized and interactive experiences for users, exhibiting the behavior of conversational AI agents \cite{gao2018neural}. However, generating instruction datasets is a crucial challenge in building instruct-based LLMs, with existing datasets ranging from crowdsourced to generated. Hand-crafted instruction instances are available in \cite{supernaturalinstructions}, while leveraging previously crowdsourced NLP datasets is a less labor-intensive curation approach \cite{wei2021flan,longpre2023flan,naturalinstructions,iyer2022opt}. LLMs have been explored for data generation in \cite{schick-schutze-2021-generating,lee2021neural,liu-etal-2022-wanli,alpaca}, and Self-Instruct \cite{wang2022self} proposes a semi-automated process for instruction instance generation. Unnatural-Instruction \cite{honovich2022unnatural} collects instruction instances by prompting a language model with only three seed examples and paraphrasing the generated instances to expand the dataset. There is also a large chunk of work that has proposed methods for automatic dataset creation \cite{li2022controllable,kim2022soda,chen2023places,meng2022generating,chen2022weakly,sahu2022data,kim2021linda,rosenbaum2022linguist,Zhang2020GroundedCG,kulhanek2021augpt,Zhang2020DialogueDO,Papangelis2021GenerativeCN,bae2020}.
Another important challenge is prompt engineering. The quality of the prompt used to guide LLMs significantly affects its performance \cite{radford2019language,brown2020language,li2021prefix}. While LMs pre-trained on large data can implicitly learn tasks with few-shot prompting, hand-crafted prompts may not always suffice. Automated prompt generation methods have been proposed, such as gradient-guided search \cite{shin2020autoprompt}, mining-based and paraphrasing-based techniques \cite{jiang2020can}, a meta-prompt \cite{reynolds2021prompt}, and automatic instruction selection and generation \cite{zhou2023large}. In this work, we introduce a conversational LLM auto-prompting method called \textit{Inception Prompting}, which enables agents to prompt each other to solve tasks through \textit{Role-Playing}. The AI user continuously provides instructions to the AI assistant for task-solving. This enables us to save the streaming instruction-solution pairs and create diverse, instructional, conversational, and task-oriented datasets. These datasets can be used to analyze the behavior and capabilities of LLMs and for future research for fine-tuning LLMs with conversational instructions.

\mysection{AI Alignment} AI alignment is a field that aims to ensure that AI systems adhere to their intended goals, interests, and values, as envisioned by their designers \cite{andreas2015alignment,hadfield2019legible,Stray2020AligningAO,Gabriel2020ArtificialIV,hadfield2021principal,Matthews2022TheAP,bai2022constitutional}. The first attempt at AI alignment was made through the "Three Laws of Robotics," which was introduced by Isaac Asimov in his science fiction stories \cite{asimov1940robot}. Developing aligned AI systems is crucial for achieving desired objectives while avoiding unintended consequences. Research in AI alignment focuses on discouraging AI models from producing false, offensive, deceptive, or manipulative information that could result in various harms \cite{kenton2021alignment,tamkin2021understanding,henderson2018ethical,goldstein2023generative}. Achieving a high level of alignment requires researchers to grapple with complex ethical, philosophical, and technical issues.
We conduct extensive experiments to study different \textit{role-playing} situations, which probe the alignment of LLMs.
